\chapter{Reproducibility Notes}
\label{app:reproducibility}

This appendix lists everything required to reproduce the empirical results
in Chapter~\ref{ch:results} from a fresh clone of the repository. The
canonical source of truth is the GitHub repository tagged at the W14 M4
freeze point; the artefacts in this thesis correspond to that tag.

\section{Repository}
\label{sec:repository}

The full source code, the evaluation scripts, the LaTeX manuscript, and the
weekly reports are released under the MIT licence at
\texttt{https://github.com/<placeholder>/reasonguard}. The annotated
evaluation dataset and the published bound batch are released on Zenodo at
\texttt{https://doi.org/<placeholder>}.

\section{Software Versions}
\label{sec:software-versions}

The canonical evaluation run uses the following versions. Versions are
captured in \texttt{requirements.txt} and pinned to ensure reproducibility
across machines.

\begin{itemize}
  \item Python 3.10+ (3.14 on the reference workstation).
  \item pandas $\geq$ 2.0, numpy $\geq$ 1.24.
  \item spaCy 3.8 with \texttt{en\_core\_web\_lg} 3.8.
  \item MLflow 2.10+, sqlite store at the URI in
        \texttt{src/pipeline\_config.MLFLOW\_TRACKING\_URI}.
  \item Ollama 0.23.3 with 4-bit quantisation on a CPU-only inference
        path.
  \item matplotlib 3.7+ for the figures.
  \item streamlit 1.30+ for the verification dashboard.
  \item pytest 7.4+ with the configuration in \texttt{pyproject.toml}.
  \item ruff for linting (lint-clean on the W14 tag).
\end{itemize}

\section{Reproduction Commands}
\label{sec:reproduction-commands}

The end-to-end pipeline runs from a clean clone in three blocks. The
script at \texttt{scripts/run\_full\_pipeline.sh} wraps these commands.

\begin{lstlisting}[language=bash,caption={Repository setup.}]
git clone https://github.com/<placeholder>/reasonguard
cd reasonguard
python -m venv .venv
source .venv/bin/activate
pip install -r requirements.txt
python -m spacy download en_core_web_lg
\end{lstlisting}

\begin{lstlisting}[language=bash,caption={Local-model setup.}]
brew install ollama          # or platform equivalent
ollama serve &                # start the local server
ollama pull llama3.1:8b-instruct-q4_0
ollama pull mistral:7b-instruct-v0.3-q4_0
ollama pull phi3:mini
\end{lstlisting}

\begin{lstlisting}[language=bash,caption={End-to-end run.}]
# One command runs the full pipeline:
scripts/run_full_pipeline.sh 25     # 25 prompts per event type

# Or the individual stages:
REASONGUARD_STRATIFY=1 REASONGUARD_PER_EVENT_TYPE=25 \
  python -m src.formal_bound_builder
REASONGUARD_STRATIFY=1 REASONGUARD_PER_EVENT_TYPE=25 \
  python -m src.prompt_builder
python -m src.synthetic_response_generator
python -m src.ollama_llm_runner
python -m src.cloud_llm_runner        # optional, needs API keys
python -m src.merge_ollama_responses
python -m src.reason_guard_checker
python -m src.event_analysis
python -m src.viz
python -m src.evaluation_metrics      # requires annotation_results.csv
python -m src.prepare_meeting_pack
python -m src.status_cli
\end{lstlisting}

\section{MLflow Experiments}
\label{sec:mlflow-experiments}

Five MLflow experiments hold the run history (declared in
\texttt{src/pipeline\_config.py}):

\begin{itemize}
  \item \texttt{reasonguard\_formal\_bounds} — one run per bound build.
  \item \texttt{reasonguard\_prompts} — one run per prompt build.
  \item \texttt{reasonguard\_synthetic\_responses} — one run per synthetic
        batch.
  \item \texttt{reasonguard\_ollama\_responses} — one run per (model, batch),
        with parameters (\texttt{model, temperature, top\_p, pilot\_limit})
        and metrics (\texttt{ok\_count, ok\_rate, mean\_elapsed\_seconds}).
  \item \texttt{reasonguard\_violation\_checker} — one run per checker
        invocation, with the V1--V5 totals as metrics.
\end{itemize}

The full run history can be inspected with:

\begin{lstlisting}[language=bash,caption={Inspect the MLflow store.}]
mlflow ui --backend-store-uri sqlite:///mlflow.db
\end{lstlisting}

\section{Test Suite}
\label{sec:test-suite}

Every module under \texttt{src/} has a corresponding test file under
\texttt{tests/}. The full suite has 80+ cases (one skip on the optional
mlflow check) and runs in under nine seconds on the reference workstation.

\begin{lstlisting}[language=bash,caption={Run the test suite.}]
python -m pytest tests/ -q
\end{lstlisting}

\section{Dataset Provenance}
\label{sec:dataset-provenance}

The IEC-104 dataset used in this thesis is the eon-iec subset of the VUT
Brno ICS Dataset for Smart Grid Anomaly Detection
\autocite{MatousekRysavyGrofcik2022SmartGrid}. The exact CSV file is
\texttt{data/raw/sample\_003.csv}, with 20\,000 atomic IEC-104 rows. The
working subset is checked into the Zenodo release; the raw dataset is
downloaded from IEEE DataPort under the DOI in the BibTeX entry.

The Modbus dataset for the function code violation and replay attack
event types is the VUT Brno Modbus Dataset for ICS Anomaly Detection
\autocite{RysavyMatousek2021Modbus}; integration of this dataset is the
W12 follow-up.
